\documentclass[11pt,a4paper,roman]{moderncv}

\moderncvstyle{banking}
\moderncvcolor{blue}
\nopagenumbers{}

\usepackage[utf8]{inputenc}
\usepackage{fontawesome5}
\usepackage{tabularx}
\usepackage{ragged2e}
\usepackage[scale=0.86]{geometry}
\usepackage{multicol}
\usepackage{import}

% personal data
\name{Ma}{Jianfa}
\title{Curriculum Vitae}

\newcommand*{\customcventry}[7][.25em]{
  \begin{tabular}{@{}l} 
    {\textbf{#4}}        % 改这里
  \end{tabular}
  \hfill
  \begin{tabular}{l@{}}
    {\textbf{#5}}        % 改这里
  \end{tabular} \\
  \begin{tabular}{@{}l} 
    {\textit{#3}}        % 改这里
  \end{tabular}
  \hfill
  \begin{tabular}{l@{}}
    {\textit{#2}}        % 改这里（如果你也想日期斜体）
  \end{tabular}
  \ifx&#7&%
  \else{\\%
    \begin{minipage}{\maincolumnwidth}%
      \small#7%
  \end{minipage}}\fi%
  \par\addvspace{#1}}

\newcommand*{\customcvproject}[4][.25em]{
  \begin{tabular}{@{}l}
    {\bfseries #2}
  \end{tabular}
  \hfill
  \begin{tabular}{l@{}}
    {\itshape #3}
  \end{tabular}
  \ifx&#4&%
  \else{\\%
    \begin{minipage}{\maincolumnwidth}%
      \small#4%
    \end{minipage}}\fi%
  \par\addvspace{#1}}

\setlength{\tabcolsep}{12pt}

\begin{document}

\makecvtitle
\vspace*{-16mm}

\begin{center}
  \renewcommand{\arraystretch}{1.3}
  \begin{tabular}{ c c }
    \faGithub\enspace \href{https://github.com/Pluviophil3}{github.com/Pluviophil3} \quad
    \faGlobe\enspace \href{https://jianfama.github.io}{jianfama.github.io} \\
    \faEnvelope\enspace \href{mailto:majianfa2004@gmail.com}{majianfa2004@gmail.com} \quad
    \faMobile\enspace (+86) 133 2280 5128\\
  \end{tabular}
\end{center}

\vspace*{-2.5mm}

%-----------------------------------------------------------------------
\section{EDUCATION}
%-----------------------------------------------------------------------

{\customcventry{Sep 2022 -- Jun 2026}{B.E. in Computer Science and Technology}{Southern University of Science and Technology}{Shenzhen, China}{}{}}
\begin{itemize}
  \item GPA: 3.6 / 4.0
  \item Centesimal grade average: 87.0
\end{itemize}
\vspace{2mm}

{\customcventry{Sep 2026 -- Jul 2027 \textnormal{\textit{(Incoming)}}}{M.Sc. in Computer Science}{The Chinese University of Hong Kong}{Hong Kong SAR, China}{}{}}
\begin{itemize}
  \item[--] Admitted; expected to enroll Fall 2026
\end{itemize}

%-----------------------------------------------------------------------
\section{RESEARCH \& VISITING EXPERIENCE}
%-----------------------------------------------------------------------

{\customcventry{Apr 2025 -- Oct 2025}{Research Intern}{SmarT Autonomous Robotics (STAR) Group, SUSTech}{Shenzhen, China}{}{}}
\begin{itemize}
  \item[--] Supervisor: Prof.\ Boyu Zhou, Dept.\ of Mechanical and Energy Engineering
  \item[--] Research on VLM-based aerial object navigation; led to the AirHunt paper (Under Review, IEEE-TASE)
\end{itemize}
\vspace{2mm}

{\customcventry{Sep 2025 -- Dec 2025}{Semester Exchange}{The Hong Kong Polytechnic University}{Hong Kong}{}{}}
\begin{itemize}
  \item[--] Program: Artificial Intelligence Engineering, Dept.\ of Electronic \& Electrical Engineering
  \item[--] Core Courses: Deep Learning, Multimedia Communication, Web System and Technology
\end{itemize}
\vspace{2mm}

{\customcventry{Jul 2024 -- Aug 2024}{Summer Workshop}{National University of Singapore}{Singapore}{}{}}
\begin{itemize}
  \item[--] Program: Visual Computing, School of Computing (SoC)
  \item[--] Supervisor: Prof.\ Terence Sim; research on traffic sign recognition in extreme scenarios
\end{itemize}
%-----------------------------------------------------------------------
\section{PUBLICATIONS}
%-----------------------------------------------------------------------

\begin{itemize}
  \item G.~Lan*, \textbf{J.~Ma*}, Z.~Zhao, Q.~Qi, Y.~Liu \& Q.~Hao.
    \textbf{RoGaussianOCC: A Gaussian Fusion Method for Accurate and Robust Occupancy Prediction in Autonomous Driving.}
    \emph{Information Fusion}. (Under Review, Jan.\ 2026) \textbf{[Co-first Author]}
  \vspace{2mm}
  \item X.~Chen, Z.~Liu, \textbf{J.~Ma}, B.~Du, T.~Zhang, X.~Wang \& B.~Zhou.
    \textbf{AirHunt: Bridging VLM Semantics and Continuous Planning for Efficient Aerial Object Navigation.}
    \emph{IEEE Transactions on Automation Science and Engineering}. (Under Review, Jan.\ 2026)
\end{itemize}

%-----------------------------------------------------------------------
\section{RESEARCH EXPERIENCE}
%-----------------------------------------------------------------------

{\customcvproject{AirHunt: Bridging VLM Semantics and Continuous Planning for Efficient Aerial Object Navigation}{Apr 2025 -- Oct 2025}
  {\begin{itemize}
    \vspace*{1mm}
    \item \textbf{Role:} Research Intern $|$ Supervisor: Prof.\ Boyu Zhou
    \vspace*{0.75mm}
    \item \textbf{Objective:} Developed a real-time aerial navigation system for large-scale unknown outdoor environments driven by natural language commands, bridging high-latency VLM reasoning with low-latency UAV control for zero-shot object navigation.
    \vspace*{0.75mm}
    \item \textbf{Method:} (1) \textit{Dual-pathway Asynchronous Architecture}: decouples slow VLM inference from fast UAV replanning; (2) \textit{Active Dual-Task Reasoning}: keyframe selection for compact scene memory; (3) \textit{Semantic-Geometric Coherent Planning}: integrates VLM semantic guidance with geometric map for continuous navigation.
    \vspace*{0.75mm}
    \item \textbf{Achievement:} 78\% navigation success rate in AirSim (98\% of human performance); completed tasks in 37\% of the flight time vs.\ existing methods; validated on custom-built quadrotor across 10$+$ hours of outdoor real-world experiments.
    \vspace*{0.75mm}
    \item \textbf{Contribution:} Designed high-frequency geometric path planner, low-frequency VLM semantic reasoning module, and ROS-based cross-module communication; implemented baselines and ablation studies; built and deployed the AirHunt quadrotor platform.
  \end{itemize}}
}

\vspace*{1.2mm}

{\customcvproject{RoGaussianOCC: Accurate and Robust Gaussian Occupancy Prediction in Autonomous Driving}{Sep 2024 -- Apr 2025}
  {\begin{itemize}
    \vspace*{1mm}
    \item \textbf{Role:} Research Leader $|$ Supervisor: Prof.\ Qi Hao \& Dr.\ Gongjin Lan
    \vspace*{0.75mm}
    \item \textbf{Objective:} Built a robust 3D Gaussian occupancy prediction model for vision-based autonomous driving, tackling random Gaussian initialization, occlusion from single-frame inference, and sparse supervision signals.
    \vspace*{0.75mm}
    \item \textbf{Method:} (1) \textit{Historical Gaussian Fusion}: temporal Gaussian propagation; (2) \textit{4D Gaussian Sampling}: aggregates historical and current frames to reduce occlusion; (3) \textit{Perspective Supervision}: auxiliary signal for faster convergence.
    \vspace*{0.75mm}
    \item \textbf{Achievement:} 32.20\% IoU and 20.65\% mIoU on nuScenes (+7.95\% / +8.12\% over baseline); validated on extreme-condition dataset; paper submitted to \emph{Information Fusion} (under review); demo and code on GitHub.
    \vspace*{0.75mm}
    \item \textbf{Contribution:} Proposed and implemented full RoGaussianOCC framework; conducted model training, parameter tuning, ablation studies, and visualization analysis; led project management; authored paper as co-first author.
  \end{itemize}}
}

\vspace*{1.2mm}

{\customcvproject{Natural Language Instructions Based 3D Cross-Modal Perception System for UAVs}{Dec 2023 -- Sep 2024}
  {\begin{itemize}
    \vspace*{1mm}
    \item \textbf{Role:} Core Member $|$ Supervisor: Prof.\ Feng Zheng
    \vspace*{0.75mm}
    \item \textbf{Objective:} Developed a generalized embodied UAV agent leveraging VLMs for natural language instruction and environmental interaction.
    \vspace*{0.75mm}
    \item \textbf{Method:} RGB-D + SLAM for 3D point cloud mapping; 3D-LLMs for cross-modal Q\&A and object detection; fine-tuned VLMs for waypoint generation; exploration path planning from discrete waypoints.
    \vspace*{0.75mm}
    \item \textbf{Contribution:} Trained and fine-tuned 3D-LLM for interaction tasks; fine-tuned VLM for waypoint generation; developed and deployed ROS system on the UAV platform.
  \end{itemize}}
}

%-----------------------------------------------------------------------
\section{COMPETITIONS \& AWARDS}
%-----------------------------------------------------------------------

\begin{minipage}{\maincolumnwidth}
  \begin{itemize}
    \item \textbf{Excellent Creativity Award (2nd Place)}, Meituan 3rd Low-Altitude Economy Intelligent Flight Management Challenge \hfill Sep 2025
    \vspace{1mm}
    \item \textbf{Third Prize}, Guangdong Contemporary Undergraduate Mathematical Contest in Modelling \hfill Dec 2023
  \end{itemize}
\end{minipage}

%-----------------------------------------------------------------------
\section{SKILLS}
%-----------------------------------------------------------------------

\begin{tabular}{ @{} >{\bfseries}l @{\hspace{4ex}} l }
  Languages           & Fluent in English (TOEFL 102), Cantonese; Native in Mandarin \\
  Programming         & Python, C/C++, Java, MATLAB, SQL \\
  Frameworks \& Tools & PyTorch, ROS, MMCV, OpenCV, TensorRT, TensorFlow, AirSim, Unreal Engine, Nvidia Jetson, Git \\
  Datasets            & nuScenes, KITTI, Waymo Open Dataset \\
  Expertise           & 3D Scene Perception \& Occupancy Prediction; Gaussian-based Representation \& Fusion; \\
                      & UAV Navigation, Motion Planning \& SLAM; VLM Fine-tuning \& Deployment; \\
                      & Embodied Agent System Development \\
\end{tabular}

\cfoot{\vspace{-2mm}{\color{gray} \rule[-10pt]{14.3cm}{0.05em}} \vspace{2mm}\\ \emph{\textcolor{gray}{Last Updated: March 2026}}}

\end{document}